\par\Needspace{6\baselineskip}\medskip
\begin{center}\hypertarget{native-174-article127}{}\pdfbookmark[2]{127.}{article-127}127.\end{center}
Sobald die drei heliozentrischen Orte auf die im vorigen Artikel beschriebene Weise gefunden
sind, kann man auch folgendermaßen vorgehen. Man bestimmt die übrigen Elemente nach dem in den
Artikeln 85–105 behandelten Verfahren zunächst aus dem ersten und dem zweiten Ort samt dem
zugehörigen Zeitintervall, dann ebenso aus dem zweiten und dem dritten Ort samt dem zugehörigen
Zeitintervall. So ergeben sich für jedes Element zwei Werte, und zwei beliebige der Differenzen
zwischen diesen Werten kann man für $X$ und $Y$ nehmen. Ein beträchtlicher Vorteil dieses Verfahrens
besteht darin, dass man bei den ersten Hypothesen alle übrigen Elemente bis auf die beiden außer
Acht lassen kann, die man zur Bestimmung von $X$ und $Y$ auswählt. Die anderen werden erst in der
letzten Rechnung auf der Grundlage der korrigierten Werte von $x,y$ bestimmt, und zwar allein aus der
ersten Kombination, allein aus der zweiten oder, was meist vorzuziehen ist, aus der Kombination des
ersten mit dem dritten Ort. Die Wahl der beiden Elemente ist im Allgemeinen beliebig und ermöglicht
viele verschiedene Lösungen. So kann man etwa den Logarithmus des Halbparameters mit dem Logarithmus
der großen Halbachse, einen dieser beiden Logarithmen mit der Exzentrizität oder eines dieser drei
Elemente mit der Länge des Perihels verbinden. Ergibt sich eine elliptische Bahn, so kann man auch
eines dieser vier Elemente mit der exzentrischen Anomalie verbinden, die in beiden Rechnungen dem
mittleren Ort entspricht. Die Formeln 27–30 in Artikel 96 machen die Rechnung dann besonders
einfach. In Sonderfällen ist bei dieser Wahl jedoch eine gewisse Umsicht nötig. Bei Bahnen, die sich
einer Parabel nähern, wären etwa die große Halbachse $a$ oder ihr Logarithmus weniger geeignet, weil
ihre übermäßig großen Änderungen kaum als proportional zu den Änderungen von $x,y$ gelten könnten. In
einem solchen Fall wäre die zweckmäßigere Wahl $\dfrac{1}{a}\mathpunct{\cdot}$ Auf diese Vorsichtsmaßnahmen gehen wir nur kurz
ein, da das fünfte Verfahren, das im nächsten Artikel erläutert wird, den vier bisher dargestellten
in fast allen Fällen überlegen ist.
